\noindent\textbf{Elección del procedimiento.} La presión manométrica es $p(h)=\rho gh$. La fuerza es la integral de esa presión, y la posición se obtiene igualando momentos: $Fh_{cp}=\int h\,p\,dA$. Para el rectángulo, $h_c=h_t+H/2$ e $I_G=bH^3/12$.

\begin{enumerate}

\item El centroide está a media altura bajo el borde superior.

\[\fitmath{\mathrm{hc}=\mathrm{ht}+\frac{H}{2}}\]

\[\fitmath{\mathrm{hc}=(1)+\frac{(3)}{2}\approx 2.5\,\mathrm{m}}\]

\item Calcular el área sobre la que actúa la presión.

\[\fitmath{A=b\cdot H}\]

\[\fitmath{A=(2)\cdot (3)\approx 6\,\mathrm{m^{2}}}\]

\item La integral de una presión lineal equivale a la presión en el centroide por el área.

\[\fitmath{F=\rho\cdot g\cdot \mathrm{hc}\cdot A}\]

\[\fitmath{F=(1000)\cdot (9.81)\cdot (2.5)\cdot (6)\approx 147150\,\mathrm{N}}\]

\item La integral del momento requiere la inercia geométrica respecto al eje centroidal horizontal.

\[\fitmath{\mathrm{IG}=\frac{b\cdot {H}^{3}}{12}}\]

\[\fitmath{\mathrm{IG}=\frac{(2)\cdot {(3)}^{3}}{12}\approx 4.5\,\mathrm{m^{4}}}\]

\item De $Fh_{cp}=\rho g(I_G+Ah_c^2)$ se despeja la profundidad de aplicación.

\[\fitmath{\mathrm{hcp}=\mathrm{hc}+\frac{\mathrm{IG}}{\mathrm{hc}\cdot A}}\]

\[\fitmath{\mathrm{hcp}=(2.5)+\frac{(4.5)}{(2.5)\cdot (6)}\approx 2.8\,\mathrm{m}}\]

\end{enumerate}
